\id{IRSTI 28.23.37}{}

\begin{header}
\swa{M.A.~Bisenbi, O.G.~Abdiramanov, M.B.~Rakhimzhanova, G.B.~Balgozhina, G.S.~Shuteyeva, D.C.~Zhamangarin}{ADAPTIVE NONLINEAR CONTROL OF SPACECRAFT ATTITUDE USING NEURAL NETWORKS}

\tsp{1}M.A.~Bisenbi\alink{https://orcid.org/0000-0003-2109-4512},
\tsp{1}O.G.~Abdiramanov\alink{https://orcid.org/0009-0004-9847-9669}\envelope,
\tsp{2}M.B.~Rakhimzhanova\alink{https://orcid.org/0000-0002-1328-8109},
\tsp{3}G.B.~Balgozhina\alink{https://orcid.org/0000-0001-9064-3330},
\tsp{2}G.S.~Shuteyeva\alink{https://orcid.org/0000-0002-0289-0874},
\tsp{1}D.C.~Zhamangarin\alink{https://orcid.org/0000-0002-2526-6492}
\end{header}

\begin{affil}
\tsp{1}Eurasian National University named by L.N. Gumilyev, Astana, Kazakhstan,

\tsp{2}Astana IT University, Astana, Kazakhstan,

\tsp{3}Abai Kazakh National Pedagogical University, Astana, Kazakhstan

\corrauthor{Corresponding-author: risbajbdiramanov@gmail.com}
\end{affil}

Recent developments in aerospace control systems emphasize the need for
intelligent and resilient algorithms capable of handling the complex
nonlinearities and uncertainties of spacecraft dynamics. This research
presents a neural adaptive fixed, time control framework to the attitude
stabilization and vibration suppression of a flexible spacecraft that is
subjected to unknown and time, varying conditions. The research object
is a spacecraft that has uncertain inertia and is equipped with flexible
appendages and is also subjected to external disturbances. The proposed
method introduces a neural network compensator into a fixed, time
backstepping controller, where the neural network approximates the
unknown nonlinear dynamics and unmodeled disturbances online. The
adaptive device in the loop empowers the controller to change its
parameters on its own, thus it is able to ensure precise attitude
tracking and vibration damping even when there is actuator saturation.

Simulation, based tests reveal that the newly designed controller is
capable of attaining quicker conver\-gence, better robustness, and lower
residual vibration as compared to conventional adaptive and sliding,
mode schemes. The findings show that neural compensation has a major
effect in stability margins and control accuracy. The advanced technique
extends the nonlinear spacecraft control domain by creating a self,
organizing structure that has the features of autonomy, adaptability,
and reliability, thus, it is a viable option for the next generation of
intelligent spacecraft guidance and stabilization systems.

{\bfseries Keywords:} adaptive control system, nonlinear dynamics,
spacecraft, stabilization, neural networks, control systems.

\begin{header}
АДАПТИВНОЕ НЕЛИНЕЙНОЕ УПРАВЛЕНИЕ ПОЛОЖЕНИЕМ КОСМИЧЕСКОГО ЛЕТАТЕЛЬНОГО АППАРАТА С ИСПОЛЬЗОВАНИЕМ НЕЙРОННЫХ СЕТЕЙ

\tsp{1}М.А.~Бейсенби,
\tsp{1}Ө.Ғ.~Әбдіраманов\envelope,
\tsp{2}М.Б.~Рахимжанова,
\tsp{3}Г.Б.~Балгожина,
\tsp{2}Г.С.~Шутеева,
\tsp{1}Д.С.~Жамангарин
\end{header}

\begin{affil}
\tsp{1}Евразийский Национальный университет имени Л.Н.Гумилева, Астана, Казахстан,

\tsp{2}Astana IT University, Астана, Казахстан,

\tsp{3}Казахский национальный педагогический университет им. Абая, Астана, Казахстан,

e-mail: risbajbdiramanov@gmail.com
\end{affil}

Последние разработки в области аэрокосмических систем управления
подчеркивают необходимость в интеллектуальных и устойчивых алгоритмах,
способных справляться со сложными нелинейностями и неопределенностями
динамики космических аппаратов. В этом исследовании рассматривается
нейронная адаптивная система управления с фиксированным временем для
стабилизации положения и подавления вибрации гибких космических
аппаратов, работающих в неизвестных и изменяющихся во времени условиях.
Объектом исследования является космический аппарат с неопределенной
инерцией и гибкими отростками, подверженными внешним воздействиям.
Предлагаемый подход интегрирует компенсатор нейронной сети в контроллер
обратного хода с фиксированным временем, где нейронная сеть
аппроксимирует неизвестную нелинейную динамику и немоделированные
возмущения в режиме реального времени. Этот адаптивный механизм
позволяет контроллеру автоматически регулировать свои параметры,
обеспечивая точное отслеживание положения и демпфирование вибрации даже
при перегрузке привода. Эксперименты, основанные на моделировании,
показывают, что предлагаемый контроллер обеспечивает более быструю
конвергенцию, повышенную надежность и снижение остаточной вибрации по
сравнению с традиционными схемами адаптивного и скользящего режимов.
Результаты подтверждают, что нейронная компенсация значительно повышает
запас устойчивости и точность управления. Разработанный метод вносит
свой вклад в область нелинейного управления космическими аппаратами,
предоставляя самоорганизующуюся структуру, которая повышает
автономность, адаптивность и надежность, предлагая перспективное
направление для будущих интеллектуальных систем наведения и стабилизации
космических аппаратов.

{\bfseries Ключевые слова:} адаптивная система управления, нелинейная
динамика, космический летательный аппарат, стабилизация, нейронные сети,
системы управления.

\begin{header}
НЕЙРОНДЫҚ ЖЕЛІЛЕРДІ ҚОЛДАНА ОТЫРЫП, ҒАРЫШТЫҚ ҰШУ АППАРАТТАРЫНЫҢ ОРНАЛАСУЫН АДАПТИВТІ СЫЗЫҚТЫҚ ЕМЕС БАСҚАРУ

\tsp{1}М.А.~Бейсенби,
\tsp{1}Ө.Ғ.~Әбдіраманов\envelope,
\tsp{2}М.Б.~Рахимжанова,
\tsp{3}Г.Б.~Балгожина,
\tsp{2}Г.С.~Шутеева,
\tsp{1}Д.С.~Жамангарин
\end{header}

\begin{affil}
\tsp{1}Л.Н.Гумилев атындағы Еуразия Ұлттық университеті, Астана, Қазақстан,

\tsp{2}Astana IT University, Астана, Қазақстан,

\tsp{3}Абай атындағы Қазақ Ұлттық Педагогикалық университеті, Астана, Қазақстан,

e-mail: risbajbdiramanov@gmail.com
\end{affil}

Аэроғарыштық басқару жүйелерінің соңғы әзірлемелері ғарыш аппараттары
динамикасының күрделі сызықтық емес және белгісіздіктерімен күресуге
қабілетті интеллектуалды және тұрақты алгоритмдердің қажеттілігін
көрсетеді. Бұл зерттеу белгісіз және уақыт бойынша өзгеретін жағдайларда
жұмыс істейтін икемді ғарыш аппараттарының орнын тұрақтандыру және
дірілін басу үшін нейрондық, бейімделетін, белгіленген уақытты басқару
жүйесін қарастырады. Зерттеу нысаны-белгісіз инерциясы бар және сыртқы
әсерлерге ұшырайтын икемді процестері бар ғарыш аппараты. Ұсынылған
тәсіл нейрондық желінің компенсаторын белгіленген уақыттағы кері
контроллерге біріктіреді, мұнда нейрондық желі нақты уақыт режимінде
белгісіз сызықтық емес динамика мен имитацияланбаған бұзылуларға
жуықтайды. Бұл адаптивті механизм контроллерге оның параметрлерін
автоматты түрде реттеуге мүмкіндік береді, тіпті диск жетегі шамадан тыс
жүктелген кезде де позицияны дәл бақылауға және дірілді демпферлеуге
мүмкіндік береді. Модельдеуге негізделген эксперименттер ұсынылған
контроллердің дәстүрлі адаптивті және жылжымалы режим схемаларымен
салыстырғанда жылдам конвергенцияны, сенімділікті арттыруды және қалдық
дірілді азайтуды қамтамасыз ететінін көрсетеді. Нәтижелер нейрондық
өтемақы тұрақтылық пен басқару дәлдігін айтарлықтай арттыратынын
растайды. Әзірленген әдіс болашақ интеллектуалды ғарыш аппараттарын
бағыттау және тұрақтандыру жүйелері үшін перспективалық бағытты ұсына
отырып, автономияны, бейімделуді және сенімділікті арттыратын өзін-өзі
ұйымдастыратын құрылымды қамтамасыз ету арқылы ғарыштық аппараттарды
сызықтық емес басқару саласына үлес қосады.

{\bfseries Түйін сөздер:} адаптивті басқару жүйесі, сызықтық емес динамика,
ғарыштық ұшу аппараты, бағдарлауды тұрақтандыру, нейрондық желілер,
басқару жүйелері.

\begin{multicols}{2}
{\bfseries Introduction.} Modern spacecraft are increasingly using large
and lightweight flexible structures such as solar panels, booms, and
antennas to meet the growing needs of long-duration and multifunctional
missions {[}1{]}. These flexible components are necessary to ensure
sufficient power supply, communications, and placement of instruments in
orbit. However, they create a strong connection between the motion of a
solid body and elastic vibrations, which significantly complicates the
dynamics of spacecraft orientation and control design. {[}2{]}. Even
small vibrations can propagate through the structure and deteriorate
pointing precision, stability, and imaging performance. Consequently,
maintaining high-precision position stabilization in conditions of
flexibility, environmental disturbances and uncertainty of parameters
remains a constant problem in aerospace engineering {[}3{]}.

The task of controlling the orientation of a spacecraft becomes even
more difficult if practical limitations are taken into account, such as
drive saturation, limited control torque, and nonlinear dynamics with
uncertain inertia parameters {[}4{]}. These factors can lead to
fluctuations, performance degradation, or even instability if
mismanaged. Over the past few decades, researchers have proposed a
variety of management strategies to address these issues. Classical
linear controllers are simple, but they cannot cope with non-linear
communication and high-angle maneuvers {[}5{]}. Reliable control methods
can provide limited stability over known ranges of disturbances, but
they often require highly accurate system models and, as a rule, are
conservative in performance {[}6{]}.

Adaptive control has been widely studied to compensate for parameter
uncertainties in real time, allowing online assessment of unknown
dynamics {[}7{]}. Although adaptive circuits can guarantee asymptotic
stability, they usually depend on constant excitation and cannot ensure
convergence in a fixed or finite time. In addition, their performance
may deteriorate with rapid changes in external torques or design
flexibility. To improve the accuracy and adaptability of control,
intelligent control methods integrating machine learning methods, in
particular neural networks (NNS), have recently attracted increasing
attention to the use of spacecraft {[}8{]}.

Neural networks have a powerful ability to approximate nonlinear
functions and unknown dynamics of systems without using explicit
mathematical models {[}9{]}. When controlling a spacecraft, NNS can
study the correspondence between input torques and orientation
characteristics, providing real-time compensation for unmodeled
disturbances, friction, and vibration effects {[}10{]}. Several studies
have shown that NN-based adaptive controllers significantly improve
reliability, vibration suppression, and overall control accuracy
compared to traditional adaptive or sliding mode approaches {[}11{]}.
However, conventional NN-based controllers are often oriented towards
asymptotic convergence, as a result of which the convergence time
depends on the initial conditions, which may not be suitable for
time-consuming tasks such as rapid orientation maneuvers or target
tracking.

In recent years, fixed-time control has become a promising approach that
guarantees convergence of the system within a given limited time,
regardless of the initial state {[}12{]}. Unlike time-limited control,
whose convergence time depends on initial conditions, fixed-time control
ensures predictability and uniformity of performance, which is critical
for flight planning and the safety of space operations. However, most
existing fixed-time systems are based on the sliding terminal control
mode, which suffers from rattling and potential features. These problems
can cause unwanted vibrations, drive wear, or even oversaturation of
regulators, especially in flexible spacecraft systems.

To overcome these limitations, this paper proposes a neural adaptive
fixed-time backstepping control framework for nonlinear spacecraft
attitude stabilization and vibration suppression {[}13{]}. The proposed
approach integrates a neural network compensator with a fixed- time
backstepping design to approximate and cancel unknown nonlinearities,
inertia variations, and external disturbances in real time. The
supporting structure ensures smooth generation of control laws,
eliminating interference while maintaining global stability. The neural
network dynamically adjusts its parameters based on tracking errors,
allowing the system to adapt autonomously to changing conditions.

Thus, the aim of this study is to develop a control scheme that combines
the reliability of adaptive control, the ability of neural networks to
learn, and the predictability of fixed-time convergence. The resulting
controller improves the accuracy of the spacecraft' s
orientation, suppresses vibration, and is resistant to uncertainties.

Consider a rigidly flexible spacecraft equipped with deployable solar
panels or antennas that provide a link between the relative motion of a
rigid body and elastic vibrations. The dynamics of spacecraft
orientation can be represented in a special orthogonal group SO(3) with
rotational motion determined by Euler' s equations:

\begin{equation*}
J\dot{\omega} + \omega \times (J\omega) = u + d + \tau_{f}
\end{equation*}

\begin{equation}
\dot{q} = Q(q)\omega
\end{equation}

Here, \(J \in R^{3 \times 3}\) is the inertia matrix of the spacecraft,
\(\omega \in R^{3}\) is the angular velocity vector, \(u \in R^{3}\) is
the control moment, \(d \in R^{3}\ \)is an unknown limited external
disturbance (e.g., solar pressure, magnetic field), and
\(\tau_{f}\ \)denotes the torque the moment caused by vibration of
flexible elements. The unit quaternion
\(q = \left\lbrack q_{0},\ q_{v}^{Τ} \right\rbrack^{Τ}\) represents
attitude, where \(q_{0}\) and \(q_{v} \in R^{3}\) denote the scalar and
vector components, respectively, satisfying
\(q_{0}^{2} + q_{v}^{Τ}q_{v} = 1\). The matrix \(Q(q)\) is the standard
quaternion transformation matrix.

Due to uncertainties in \(J\), unmodeled dynamics, and unknown
disturbances, the exact model is not available for direct control
design. To simplify the representation, the attitude dynamics can be
written in compact nonlinear form as:

\begin{equation}
\dot{x} = f(x) + g(x)u + \Delta (x,t)
\end{equation}

where \(x = \left\lbrack q_{v}^{Τ},\ \omega^{Τ} \right\rbrack^{Τ}\) is
the system state vector, \(f(x)\) represents the nominal nonlinear
dynamics, \(g(x)\) is the control input matrix, and
\(\Delta (x,t)\) denotes the lumped uncertainty, including
modeling errors, vibration dynamics, and external disturbances.

The main control objective is to design a control input \(u(t)\) such
that the attitude \(q\) tracks a desired orientation \(q_{d}\) while
suppressing vibration, ensuring convergence of both the attitude error
\(e_{q} = q\bigotimes q_{d}^{- 1}\) and angular velocity error
\(e_{\omega} = \omega - \omega_{d}\) to a small neighborhood around zero
within a fixed time \(T_{\max}\), regardless of initial conditions.

To handle nonlinear and uncertain terms in \(f(x)\) and
\(\Delta (x,t)\), a neural network (NN) approximator is
introduced to estimate the unknown dynamics online. The NN is defined
as:

\begin{equation}
\widehat{f}(x) = W^{Τ}\varphi(x)
\end{equation}

where \(\varphi(x) \in R^{m}\) is a vector of activation functions, and
\(W \in R^{m}\) is the configurable weight vector. The approximation
error is:

\begin{equation}
\varepsilon(x) = f(x) - \widehat{f}(x)
\end{equation}

which is supposed to be limited by the known constant
\(\left\| \varepsilon(x) \right\| \leq \varepsilon_{\max}\). The law of
adaptive updating for NN weights will be developed later to minimize
this error during operation, allowing the controller to study the
nonlinear behavior of the system in real time.

This study uses a fixed-return approach, in which the laws of virtual
control are recursively developed to stabilize each subsystem. A
non-singular finite sliding surface guarantees a finite rate of
convergence, while neural compensation provides resilience to
concentrated uncertainties. By integrating the principles of adaptive
learning and fixed-time convergence, the proposed control scheme
provides fast, stable and accurate position adjustment for both rigid
and flexible communication.

{\bfseries Materials and methods.} The main goal of this study is to
develop an adaptive nonlinear control system for spacecraft position
stabilization that effectively compensates for model uncertainties and
external disturbances using neural networks. The proposed approach is
based on previous developments in adaptive control, neural network
approximation, and fixed-time convergence control {[}14{]}.

A. Dynamics and kinematics of spacecraft.

The dynamics of the orientation of a rigidly flexible spacecraft is
determined by the Euler rotation equations:

\begin{equation}
J\dot{\omega} + \omega \times (J\omega) = u + d + \tau_{f}
\end{equation}

where \(J \in R^{3 \times 3}\) is the inertia matrix of the spacecraft,
\(\omega \in R^{3}\) is the angular velocity vector, \(u \in R^{3}\) is
the control moment, \(d \in R^{3}\) are external disturbances, and
\(\tau_{f}\) is the moment caused by vibration of flexible elements.

The orientation kinematics, represented by the unit quaternions
\(q = \left\lbrack q_{0},q_{v}^{Τ} \right\rbrack^{Τ}\), is expressed as:

\begin{equation}
\dot{q} = \left( \frac{1}{2} \right)Q(q)\omega
\end{equation}

where \(Q(q)\) is the standard quaternion transformation matrix
{[}15{]}. Quaternions are used to avoid the singularities of Euler
angles and to provide a globally correct representation of the
orientation of the spacecraft {[}16{]}.

To formulate a compact nonlinear model, the full dynamics can be
rewritten as:

\begin{equation}
\dot{x} = f(x) + g(x)u + \Delta (x,\ t)
\end{equation}

where \(x = \left\lbrack q_{v}^{T},\omega^{Τ} \right\rbrack^{Τ}\) is the
vector of the system state, \(f(x)\) denotes nominal dynamics, \(g(x)\)
is the matrix of control input data, and \(\Delta (x,\ t)\)
represents concentrated uncertainty, including unmodeled flexibility,
parameter variations, and external disturbances {[}17{]}.

Neural Network - based approximation. To compensate for unknown
nonlinearities \(\Delta (x,t)\), a Radial Basis Function Neural
Network (RBFNN) is employed as a universal approximator {[}18{]}:

\begin{equation}
\widehat{\Delta }(x) = W^{Τ}\varphi(x)
\end{equation}

where \(\varphi(x) \in R^{m}\) is a vector of Gaussian basis functions,
\(W \in R^{m \times 3}\) is the ideal weight matrix, and the
approximation error
\(\varepsilon(x) = \Delta (x,\ t) - \widehat{\Delta }(x)\)
is assumed to be bounded by
\(\left\| \varepsilon(x) \right\| \leq \varepsilon_{\max}\) {[}19{]}.

The adaptive weight update law, derived from Lyapunov stability theory,
is given by

\begin{equation}
\widehat{W} = \Gamma\varphi(x)s^{Τ}
\end{equation}

where \(\widehat{W}\) is the estimate of \(W,\Gamma > 0\) is a positive-
definite learning-rate matrix, and \(s = \omega - \omega_{d}\) is the
sliding surface representing angular-velocity tracking error {[}20{]}.
This law of updating makes it possible to study dynamic uncertainties
online, while maintaining the stability of the system {[}21{]}.

B. Nonlinear control scheme. The proposed control scheme combines neural
adaptive compensation, step-by-step design, and fixed-time convergence
concepts. The backstepping system guarantees smooth generation of
virtual control and avoids failures typical of controllers with smooth
operation {[}22{]}, while the fixed-time structure ensures convergence
over a specified time, regardless of initial conditions.

The total control moment is calculated as

\begin{equation}
u = J{\dot{\omega}}_{d} + \omega \times (J\omega) - Ks - {\widehat{W}}^{T}\phi(x)
\end{equation}

where \(K > 0\) is a positive-definite gain matrix determining the rate
of convergence, and the neural- network term
\({\widehat{W}}^{Τ}\varphi(x)\) adaptively compensates for concentrated
uncertainties and external disturbances {[}23{]}.

Using the corresponding Lyapunov analysis, it is demonstrated that all
closed-loop signals remain limited, and the errors of orientation and
angular velocity converge to small controlled neighborhoods close to
zero for a fixed time {[}24{]}.

C. Implementation algorithm. The algorithm is implemented in an onboard
discrete-time environment suitable for spacecraft computers, in
accordance with the methodology described in {[}25{]}:

1. Initialization: Set the initial neural weights \(\widehat{W}(0)\) and
the parameters \(J,\ \Gamma,\) and \(K\).

2. Measurement: Get the current position \(q\) and angular velocity
\(\omega\) from sensors.

3. Calculation of the sliding surface: Calculate
\(s = \omega - \omega_{d}\).

4. Neural-weight update: Adjust \(\widehat{W}\) using Eq. (9).

5. Torque estimation: Calculate \(u\) using Eq. (10) and apply
constraints to the drive.

6. Command execution: Send control torque to actuators until
\(s \rightarrow 0\) or mission objectives are achieved.

Overload protection and filtering mechanisms are built in to maintain
stability in the presence of drive saturation and measurement noise
{[}26{]}.

D. Selection and justification of the method. Classical control methods
such as PID and LQR are simple to implement, but they are not able to
effectively cope with strong nonlinearities, flexibility, and uncertain
dynamics {[}27{]}. Reliable \emph{H∞} controllers and controllers based
on perturbation observation provide limited performance, but require
accurate modeling and can be overly conservative. Controllers operating
in sliding mode provide reliability, but at the same time suffer from
rattling and intermittent control signals.

The proposed neural adaptive reverse step method with a fixed execution
time has several advantages:

- neural networks provide real-time approximation of unmodeled
nonlinearities without explicit parameter identification.

- the reverse stroke ensures smooth and continuous operation, avoiding
rattling.

- fixed-time control guarantees predictable convergence regardless of
initial conditions, which is essential for time-consuming tasks.

This combination provides superior reliability, faster stabilization,
and improved vibration suppression compared to classical adaptive or
sliding mode approaches, making it highly efficient for nonlinear space
systems.

{\bfseries Results and discussion.} To evaluate the performance of the
proposed neural adaptive feedback controller with a fixed action time, a
number of detailed numerical calculations were performed on a nonlinear
flexible spacecraft model. The spacecraft was tested under conditions of
uncertain inertia parameters, external torque disturbances, and drive
saturation to test the effectiveness of the control circuit in terms of
position stabilization, vibration suppression, and reliability.

Configuring the simulation.

The spacecraft model consisted of a rigid central body with two flexible
solar panels. The nominal inertia matrix was defined as:

\begin{equation*}
J = \text{diag} (12,5; 11,3; 9,8) \text{ kg} \cdot \text{m}^2
\end{equation*}

and it varied by ±25\% to represent the error. The maximum moment of
disturbance was 0.5 N·m per axis, and output torque of the drives was
limited to ±1 N·m.

The radial basis function neural network (RBFNN) used 15 Gaussian nodes
per axis. The learning rate and gain matrices were tuned as:

\begin{equation*}
\Gamma = \text{diag} (10, 10, 10) \text{ and } K = \text{diag} (4,5; 4,0; 3,8)
\end{equation*}

These parameters ensured stable and fast convergence. A conventional
adaptive controller (CAC) was used as a reference under identical
conditions.

Attitude stabilization result

\fig[\columnwidth]{i/image73}[Fig.1 - Convergence of the orientation error]

Figure 1 shows the error rate of orientation determination (‖qᵥ‖) over
time. The proposed controller achieved convergence in 8.5 seconds, while
the traditional adaptive method took 22 seconds to achieve similar
accuracy. The error of the steady-state position remained below
0.01°, which confirmed convergence over a fixed time,
independent of the initial conditions.

Flexible vibration suppression

\fig[\columnwidth]{i/image74}[Fig.2 - Flexible battery tip deflection response]

The reaction of a flexible solar cell to the tip deflection is shown in
Figure 2. When using a conventional controller, the oscillations lasted
for more than 25 seconds with a residual amplitude of about 0.15 m. The
proposed controller dampened these fluctuations for 8 seconds, reducing
the deviation to less than 0.03 m.

Interference elimination and robustness.

To assess the reliability along the Y-axis, a disturbing moment of 0.5
N·m was applied at t = 15 seconds. The proposed controller restored
stability in 2 seconds, while the traditional controller took 10
seconds. Neural adaptation effectively compensated for this disruption
without hesitation or instability.

Computational performance. A comparative evaluation of both control
systems showed that the fixed-time adaptive neural controller
consistently outperforms the traditional adaptive approach in all key
performance indicators. The proposed controller provided faster and more
predictable orientation while reducing installation time by almost 60\%
and more accurate angular velocity tracking.
\end{multicols}

\fig{i/image75}[Fig.3 - Timeline chart]

\begin{multicols}{2}
Vibration suppression has been significantly improved, the residual
deflection has been reduced from 0.15 m to 0.03 m, and the time to
achieve a stable vibration level has been reduced by more than 70\%. The
control signals generated by the neural controller were smoother, with
less overspeed, and required less peak torque.

During reliability tests, the controller effectively suppressed external
interference and maintained stable operation when the drive was
overloaded. The position deviation after interference was reduced from
0.052 rad (at CAC) to 0.018 rad, which confirms the high resistance to
interference.

Despite the higher computational requirements, the data processing time
by the neural controller remained significantly below the threshold
value of the on-board system. This confirms its feasibility for
real-time spacecraft applications.

The results confirm that the proposed control strategy provides
significant advantages in speed, accuracy, vibration damping, and
reliability, all while maintaining computational efficiency. The
demonstrated performance improvements make it an excellent candidate for
autonomous spacecraft attitude control, especially for missions
involving flexible structures or large-scale deployments where
traditional controllers often do not provide consistent results.

The obtained simulation results convincingly indicate that the proposed
neural adaptive feedback controller with a fixed action time
significantly improves the characteristics of spacecraft position
stabilization in nonlinear, uncertain and time-varying conditions. The
fixed-time property ensures that the orientation and angular velocity
errors converge within a strictly limited time, regardless of the
initial state. This feature is especially useful for space missions that
require urgent execution, such as autonomous rendezvous, docking
maneuvers, debris avoidance, or rapid reorientation required during
payload operations. Compared to classic controllers with finite
operating time or sliding mode, the fixed operating time system
eliminates dependence on initial conditions and avoids the
high-frequency interference commonly seen in intermittent control laws.

The inclusion of the radial basis function neural network (RBFNN) plays
a crucial role in increasing the adaptability and reliability of the
system. By approximating unknown nonlinearities and perturbations in
real time, the neural compensator allows the controller to maintain
stability even under conditions of sudden environmental changes,
uncertain inertia, or sudden vibration effects. This real-time
adaptation mechanism allows the spacecraft to study the changing
dynamics of the system autonomously and compensate for it, improving
both transient and steady-state characteristics.

An important practical advantage of this approach is the smooth
generation of torque and predictable behavior in transients, which are
necessary for high-precision tasks such as Earth observation, satellite
formation, and payload guidance for communications. The simulation
results showed that the proposed controller provides high guidance
accuracy with a reduced load on the drive, which means reduced energy
consumption and increased drive life. Such efficiency is crucial for
small satellites and nanosatellites with limited power consumption and
hardware limitations.

However, for practical implementation, several engineering
considerations must be taken into account. Real spacecraft systems face
sensor interference, time delays, signal quantization, and limited
onboard processing capabilities. To ensure reliable operation, future
developments should use noise filtering, delay-compensated adaptation
laws, and a computationally optimized neural architecture. Another
direction involves the implementation of the controller on
hardware-in-the-loop (HIL) test benches or satellite simulators in real
time to test its performance in real conditions, including communication
latency and drive saturation dynamics.

The proposed control concept can also be extended to joint
multi-satellite systems, where convergence at a fixed time will ensure
synchronized maneuvers of several spacecraft operating in service. In
addition, combining neural adaptation with reinforcement learning or
deep actor-critic architecture can further enhance the
system' s decision-making ability by providing fully
autonomous orientation and orbit control.

The research contributes to the creation of a reliable, adaptive, and
computationally feasible framework for intelligent spacecraft control.
It bridges the gap between classical analytical control theory and
modern data-driven learning, providing the foundation for the next
generation of autonomous, self-organizing, and fault-tolerant space
systems.

{\bfseries Conclusion.} This article presents a neural adaptive feedback
controller with a fixed action time to stabilize the position of the
spacecraft and suppress vibrations in uncertain and dynamic conditions.
The framework combines a neural network with a radial basis function
(RBFNN) to compensate for unknown nonlinearities in real time with a
fixed-time design that ensures convergence within a guaranteed boundary,
regardless of the initial conditions.

The simulation results showed that the proposed controller provides
faster convergence, higher guidance accuracy and higher reliability
compared to traditional adaptive methods. The neural component improved
interference suppression and adaptability, while the feedback structure
ensured a smooth change in control moments without rattling or
instability.

This approach is suitable for missions involving flexible or deployable
structures, autonomous docking, and Earth observation, where precision
and time predictability are critical. The fixed-time property ensures
reliability when performing time-consuming operations, while the
adaptive neural layer supports long-term autonomy in the uncertain space
environment.

Future work should focus on in-loop hardware verification, optimizing
neural network architecture to minimize computation, and integration
with deep learning or reinforcement learning to enhance self-learning
capabilities.

Thus, the proposed controller represents a practical and intelligent
solution for next-generation spacecraft, combining theoretical stability
with the ability to adapt in real time to provide fast, reliable and
autonomous attitude control.
\end{multicols}

\begin{center}
{\bfseries References}
\end{center}

\begin{refs}
1. Yao Q. et al. Neural adaptive fixed-time attitude stabilization and
vibration suppression of flexible spacecraft
//Mathematics.-2022.Vol.10(10):1667.\href{https://doi.org/10.3390/math10101667}{DOI
10.3390/math10101667}.

2. Chen R., Wang Z., Che W. Adaptive sliding mode attitude-tracking
control of spacecraft with prescribed time performance //Mathematics.
-2022. -Vol.10(3):40. DOI10.3390/math10030401.

3. Zhang G. Et al. Fixed-time sliding mode attitude control of a flexible
spacecraft with rotating appendages connected by magnetic bearing
//Mathematical Biosciences and Engineering. -2022.-Vol.19(3): 2286-2309.
DOI 10.3934/mbe.2022106.

4. Javaid U., Dong H. Disturbance-observer-based attitude control under
input nonlinearity //Transactions of the Institute of Measurement and
Control. -2021.-Vol.43(10): 2358-2367. DOI 10.1177/0142331221997204.

5. Liu Y. et al. An overview of finite/fixed-time control and its
application in engineering systems//IEEE/CAA Journal of Automatica
Sinica. -2022.-Vol.9(12):2106-2120. DOI 10.1109/JAS.2022.105413.

6. Li S., Liu K., Liu M. Adaptive Dynamic Programming-Based Spacecraft
Attitude Control Under a Tube-Based Framework//Electronics.
-2024.-Vol.13(22): 4575. DOI 10.3390/electronics13224575.

7. Wu H., Xu D. Attitude regulation of flexible spacecrafts with unknown
control directions and input disturbances //Journal of Systems Science
and Complexity.-2024.-Vol.37.-P.2347-2367. DOI
10.1007/s11424-024-3539-8.

8. Lu M. et al. RBFNN-Based Adaptive Fixed-Time Sliding Mode Tracking
Control for Coaxial Hybrid Aerial--Underwater Vehicles Under
Multivariant Ocean Disturbances //Drones. -2024.-Vol.8(12): 745. DOI
10.3390/drones8120745.

9. Esmaeilzadeh S. M., Zeyghami M. S. Nonlinear finite time attitude
control of flexible spacecraft based on a novel output redefinition
method //Chinese Journal of Aeronautics. -2023.-Vol.36(11). -P.373-385.
DOI 10.1016/j.cja.2023.08.001.

10. Liu F., Yao X., Liu Y. Adaptive neural network control of
multiple-sectioned flexible riser with time-varying output constraint
and input nonlinearity //Asian Journal of Control. -2023.-Vol.26(2).
-P.917-930.DOI 10.1002/asjc.3231.

11. Efimov D. Et al. Finite-time stability tools for control and
estimation //Foundations and Trends in Systems and Control. -2020.
Vol.9(2-3): 171-364. DOI 10.1561/2600000026.

12. Post M. A., Yan X. T., Letier P. Modularity for the future in space
robotics: A review //Acta Astronautica. -2021.-Vol.189.-P.530-547. DOI
10.1016/j.actaastro.2021.09.007.

13. Liu Z. Et al. Dual-condition event-triggered distributed adaptive
neural control for multi-agent systems with nonlinear fault //Journal of
the Franklin Institute. -2022.-Vol.359(14). -P.7393-7414. DOI
10.1016/j.jfranklin.2022.07.026.

14. You Z. Et al. Adaptive neural network vibration suppression control
of flexible joints space manipulator based on H∞ theory//Journal of
Vibroengineering. -2023.-Vol.25(3). -P.492-505. DOI
10.21595/jve.2022.22797.

15. Cao T., Gong H., Han B. Observer-Based Predefined-Time Attitude
Control for Spacecraft Subject to Loss of Actuator
Effectiveness//Processes. -2022.-Vol.10(11): 2294. DOI
10.3390/pr10112294.

16. Matessa M. Algorithms for encoding interstellar messages//Acta
Astronautica. -2025. -Vol.229: 161-165. DOI
10.1016/j.actaastro.2025.01.011.

17. Czornik A. et al. Some conditions for the proportional local
assignability of the Lyapunov spectrum of discrete linear time-varying
systems //Automatica. -2022.-Vol.144: 110458. DOI
10.1016/j.automatica.2022.110458.

18. Chen W., Hu Q. Sliding-mode-based attitude tracking control of
spacecraft under reaction wheel uncertainties //IEEE/CAA Journal of
Automatica Sinica. -2023.-Vol.10(6): 1475-1487. DOI
10.1109/JAS.2022.105665.

19. Kong X. Et al. Predefined-Time Control of a Spacecraft Attitude with
Thrust Booms //Aerospace. -2023.-Vol.10(2): 94. DOI
10.3390/aerospace10020094.

20. Bima A. I. H. et al. Integrative system biology and mathematical
modeling of genetic networks identifies shared biomarkers for obesity
and diabetes //Mathematical Biosciences and Engineering.
-2022.-Vol.19(3): 2310-2329. DOI 10.3934/mbe.2022107.

21. Liu M. Et al. Adaptive dynamic programming-based fault-tolerant
attitude control for flexible spacecraft with limited wireless
resources//Science China Information Sciences. -2023.-Vol.66: 202201.
DOI 10.1007/s11432-022-3732-9.

22. Gao Y. Et al. Fault deviation estimation and integral sliding mode
control design for Lipschitz nonlinear systems //Systems \& Control
Letters. -2019.-Vol.123.-P.8-15. DOI 10.1016/j.sysconle.2018.08.006.

23. Guo Q. Et al. Neural adaptive backstepping control of a robotic
manipulator with prescribed performance constraint //IEEE Transactions
on Neural Networks and Learning Systems. -2019.-Vol.30(12).
-P.3572-3583. DOI 10.1109/TNNLS.2018.2854699.

24. Lunghi P. Et al. Energy efficiency analysis of Spiking Neural
Networks for space applications//Astrodynamics. -2025.-Vol.9(6).
-P.909-932. DOI 10.1007/s42064-024-0256-y.

25. Dai X. Et al. Mechanical behaviors of inner and outer sidewalls of
honeycomb cores subjected to out-of-plane compression //Aerospace
Science and Technology. -2022.-Vol.127:107659. DOI
10.1016/j.ast.2022.107659.

26. Abdiramanov O. Et al. Self-Organizing Control Systems for Nonlinear
Spacecraft in the Class of Structurally Stable Mappings //International
Journal of Advanced Computer Science and Applications.
-2023.-Vol.14(11). -Р.786-792. DOI 10.14569/ijacsa.2023.0141179.

27. Beisenbi M. Et al. Control of deterministic chaotic modes in power
systems //IEEE. -2023.DOI 10.1109/SIST58284.2023.10223493.
\end{refs}

\begin{info}
\hspace{1em}\emph{{\bfseries Information about authors}}

Beisenbi M.A. - Doctor of Technical Sciences, Professor, L.N.Gumilyov
Eurasian National University, Astana, Kazakhstan, e-mail:
beisenbi@mail.ru;

Abdiramanov O.G. - PhD student, L.N.Gumilyov Eurasian National
University, Astana, Kazakhstan, e-mail:
risbajbdiramanov@gmail.com;

Rakhimzhanova M.B. - PhD, Assistant Professor Astana IT University,
Astana, Kazakhstan, е-mail: rakhimzhan.mira@gmail.com;

Balgozhina G.B. - doctoral student, Abai Kazakh National Pedagogical
University, Almaty Kazakhstan, е-mail: balgulmira300982@gmail.com;

Shuteyeva G.S. - PhD, Astana IT University, Astana, Kazakhstan, е-mail:
g.shuteyeva@astanait.edu.kz

Zhamangarin D.S. - PhD, associate professor, L.N. Gumilyov Eurasian
National University, Astana, Kazakhstan, е-mail:
enu\_1@enu.kz.

\hspace{1em}\emph{{\bfseries Сведения об авторах}}

Бейсенби М.А. - д.т.н., профессор, Евразийский Национальный университет
им. Л.Н.Гумилёва, Астана, Казахстан, е-mail:
beisenbi@mail.ru;

Әбдіраманов Ө.Ғ. - докторант, Евразийский Национальный
университет им. Л.Н.Гумилёва, Астана, Казахстан, е-mail:
risbajbdiramanov@gmail.com;

Рахимжанова М.Б. - PhD, асс.профессор, Astana IT University, Астана,
Казахстан, е-mail: rakhimzhan.mira@gmail.com;

Балгожина Г.Б. - докторант, Казахский национальный педагогический
университет им. Абая, Алматы, Казахстан, е-mail:
balgulmira300982@gmail.com;

Шутеева Г.С. - PhD, Astana IT University, Астана, Казахстан, е-mail:
g.shuteyeva@astanait.edu.kz;

Жамангарин Д.С. - PhD, асс. профессор, Евразийский Национальный
университет им. Л.Н.Гумилёва, Астана, Казахстан, е-mail: enu\_1@enu.kz.
\end{info}
